\section{Introduction}

Flash loans have become an important primitive in decentralized finance because they allow users to access large amounts of liquidity without upfront collateral, as long as the borrowed assets are returned within a single blockchain transaction. This feature greatly improves capital efficiency and supports common DeFi activities such as arbitrage, liquidation, and collateral rebalancing. At the same time, it also creates new security risks because attackers can temporarily control large amounts of funds and combine multiple protocols within one atomic transaction.

The goal of this project is to understand how flash loans work in practice, how flash loan transactions can be identified from on-chain data, and how such transactions are used in the broader DeFi ecosystem. To achieve this goal, we first study the interaction process between flash loan providers and users. We then design protocol-specific detection patterns and apply them to Ethereum transaction data. Finally, we analyze the detected transactions and summarize their major characteristics and applications.
